% =====================================================================
% ri_solstyle.tex: worked solutions in the layout of the RI 2026 Year 6 Prelim
% "Solutions with comments" (grail.moe 26612/26613), loaded AFTER prelim_style.
%   * page 1: bold school-style block + bold underlined title; later pages: running
%     title top-left with a rule under it; footer "Page N of M" bottom right (Calibri).
%   * each question: prelim-format question text (bold number, marks right, source tag
%     bottom right), then a bordered table:  Solution | Comments  (bold header row);
%     body rows  label + [marks] (bold, narrow left column) | working | comments,
%     a full-width rule between parts, working left-aligned (fleqn).
% Rows come from the existing sources: \pt in qparts (Ch10 HW), \slabel (Ch10 Test),
% spart (Ch11). Marks under each label are taken from the question's \mk.
% =====================================================================
\makeatletter
\usepackage{lastpage}
\usepackage{adjustbox}
\usepackage{zref-savepos,zref-abspage}
\zref@addprop{savepos}{abspage}
\geometry{a4paper, top=22mm, bottom=24mm, left=22mm, right=22mm, headsep=8mm, footskip=12mm, headheight=24pt}
\providecommand{\risoltitle}{\prelimtitle: Solutions with comments}
\providecommand{\risolschool}{Compiled by Maximus Pung Jun}
\fancypagestyle{risol}{\fancyhf{}%
  \lhead{\fontsize{11}{12}\selectfont 2026 \risoltitle\\[-1.2ex]\rule{\textwidth}{0.6pt}}%
  \rfoot{\calibrifont\fontsize{11}{13}\selectfont Page \textbf{\thepage} of \textbf{\pageref*{LastPage}}}%
  \renewcommand{\headrulewidth}{0pt}\renewcommand{\footrulewidth}{0pt}}
\fancypagestyle{risolfirst}{\fancyhf{}%
  \rfoot{\calibrifont\fontsize{11}{13}\selectfont Page \textbf{\thepage} of \textbf{\pageref*{LastPage}}}%
  \renewcommand{\headrulewidth}{0pt}\renewcommand{\footrulewidth}{0pt}}
\pagestyle{risol}
\newif\ifri@fresh
\newcommand{\risolfront}{\thispagestyle{risolfirst}\global\ri@freshtrue%
  {\noindent\bfseries\linespread{1}\selectfont \risolschool\\ H2 Mathematics (9758)\\ 2026\par}%
  \vspace{4ex}{\noindent\bfseries\underline{2026 \risoltitle}\par}\vspace{3ex}}
% ---- table geometry (RI: label 42pt | working 284pt | comments 142pt of 468pt) ----
\newlength{\ri@L}\setlength{\ri@L}{15mm}
\newlength{\ri@C}\setlength{\ri@C}{40mm}
\newlength{\ri@pad}\setlength{\ri@pad}{1.8mm}
\newif\ifri@sol \newif\ifri@first
\newcount\ri@depth
\def\ri@outer{}
\newtcolorbox{risol}{enhanced, breakable, sharp corners, frame hidden, colback=white,
  boxrule=0pt, boxsep=0pt, left=\dimexpr\ri@L+\ri@pad\relax, right=\dimexpr\ri@C+\ri@pad\relax,
  top=1.2mm, bottom=2.5mm, toptitle=0.6mm, bottomtitle=0.6mm, lefttitle=\ri@pad, righttitle=0pt,
  colbacktitle=white, coltitle=black, fonttitle=\bfseries,
  title={Solution\hfill\makebox[\ri@C][l]{\hspace{\ri@pad}Comments}},
  before skip=2ex, after skip=2ex,
  overlay unbroken and first={\draw[line width=0.5pt] (frame.south west) rectangle (frame.north east);
     \draw[line width=0.5pt] (title.south west) -- (title.south east);
     \draw[line width=0.5pt] ([xshift=\ri@L]interior.north west) -- ([xshift=\ri@L]frame.south west);
     \draw[line width=0.5pt] ([xshift=-\ri@C]frame.north east) -- ([xshift=-\ri@C]frame.south east);},
  overlay middle and last={\draw[line width=0.5pt] (frame.south west) rectangle (frame.north east);
     \draw[line width=0.5pt] ([xshift=\ri@L]frame.north west) -- ([xshift=\ri@L]frame.south west);
     \draw[line width=0.5pt] ([xshift=-\ri@C]frame.north east) -- ([xshift=-\ri@C]frame.south east);},
  after upper={\ri@endrow\protected@write\@auxout{}{\string\ri@settotal{\ri@qid}{\the\ri@sumH}}},
  before upper={\global\ri@sumH=14mm\ri@soltrue\ri@firsttrue\ri@depth=0 \parindent=0pt \parskip=0.6ex
     \predisplaypenalty=100 
     \linespread{1.12}\selectfont\raggedright\@fleqntrue\@mathmargin=0pt
     \abovedisplayskip=5pt plus 2pt\belowdisplayskip=5pt plus 2pt
     \abovedisplayshortskip=3pt plus 1pt\belowdisplayshortskip=4pt plus 2pt
     \ifcsname ri@hp@\ri@qid\endcsname\else\ri@row{}\fi}}
\def\ri@qid{}
% one row: full-width rule (glue leaders, so it vanishes at a page break where the frame
% edge takes its place; not before the first row), then bold label + bold [marks]
% in the left column, on the baseline of the first line of the working
\ExplSyntaxOn
\tl_new:N \g__ri_tmp_tl
\cs_new_protected:Npn \ri_setkey:Nn #1#2
  { \group_begin: \cs_set_eq:NN \chg \use_ii:nn \cs_set_eq:NN \chgm \use_ii:nn
    \cs_set_eq:NN \textbf \use:n \cs_set_eq:NN \mathrm \use:n
    \tl_gset:Nx \g__ri_tmp_tl {#2} \group_end:
    \tl_set_eq:NN #1 \g__ri_tmp_tl \tl_remove_all:Nn #1 {~} }
\cs_set_eq:NN \ri@setkey \ri_setkey:Nn
\ExplSyntaxOff
% ---- page-break rules (positions from the previous run, zref labels in the .aux) ----
% * a row (one part) never breaks across pages: if it does not fit on the rest of the
%   page it starts the next page (forced break between rows); only a row taller than a
%   whole page may break inside;
% * a whole solution that fits on one page is kept together (needspace before the table);
%   a longer one starts where at least its first row fits;
% * a row is at least as tall as its label + [marks].
\newcount\ri@n \newbox\ri@labbox \def\ri@open{0}
\newdimen\ri@ybot \newdimen\ri@ytop \newdimen\ri@h \newdimen\ri@sumH
\AtBeginDocument{\ri@ybot=\dimexpr\paperheight-1in-\voffset-\topmargin-\headheight-\headsep-\textheight\relax
  \ri@ytop=\dimexpr\ri@ybot+\textheight\relax}
\def\ri@pg#1{\zref@extractdefault{#1}{abspage}{-1}}
% \ri@rowh{n}: height of row n in the previous run -> \ri@h (0pt if unknown)
\newcommand{\ri@rowh}[1]{\ri@h=0pt
  \ifnum\ri@pg{ri-s#1}>0 \ifnum\ri@pg{ri-e#1}>0
    \ifnum\ri@pg{ri-s#1}=\ri@pg{ri-e#1}\relax \ri@h=\dimexpr\zposy{ri-s#1}sp-\zposy{ri-e#1}sp\relax
    \else \ri@h=\dimexpr\zposy{ri-s#1}sp-\ri@ybot+\ri@ytop-\zposy{ri-e#1}sp-6mm\relax\fi\fi\fi}
\newcommand{\ri@setfirst}[2]{\expandafter\gdef\csname ri@first@#1\endcsname{#2}}
\newcommand{\ri@setkept}[1]{\expandafter\gdef\csname ri@kept@#1\endcsname{}}
\newcommand{\ri@settotal}[2]{\expandafter\gdef\csname ri@total@#1\endcsname{#2}}
\newcommand{\ri@endrow}{\ifnum\ri@open=1 \par\zsavepos{ri-e\the\ri@n}%
  \ifnum\ri@pg{ri-s\the\ri@n}=\ri@pg{ri-e\the\ri@n}\relax
    \dimen@=\dimexpr\ri@curH sp+1.6mm-\zposy{ri-s\the\ri@n}sp+\zposy{ri-e\the\ri@n}sp\relax
    \ifdim\dimen@>0pt \nobreak\kern\dimen@\penalty\z@\fi\fi
  \ri@rowh{\the\ri@n}\global\advance\ri@sumH\dimexpr\ri@h+3.5mm\relax
  \gdef\ri@open{0}\fi}
\newcommand{\ri@row}[1]{\par\ri@endrow\global\advance\ri@n\@ne\gdef\ri@open{1}\@beginparpenalty\@M
  \ifri@first\global\ri@firstfalse\ri@firstfalse
    \protected@write\@auxout{}{\string\ri@setfirst{\ri@qid}{\the\ri@n}}%
  \else
    \zsavepos{ri-p\the\ri@n}%
    % push the row to the next page if it would not fit on the rest of this one
    \ri@rowh{\the\ri@n}%
    \ifnum\ri@pg{ri-p\the\ri@n}>0 \ifdim\ri@h>0pt
      \ifdim\dimexpr\ri@h+4mm\relax<\dimexpr\textheight-10mm\relax
        \ifdim\dimexpr\ri@h+7mm\relax>\dimexpr\zposy{ri-p\the\ri@n}sp-\ri@ybot\relax
          \penalty-\@M\fi\fi\fi\fi
    \vskip0.6ex\nointerlineskip
    \leaders\hbox{\hskip-\dimexpr\ri@L+\ri@pad\relax\vrule width\dimexpr\linewidth+\ri@L+\ri@C+2\ri@pad\relax height0.5pt depth0pt}\vskip0.5pt
    \nobreak\vskip1.4mm\fi
  \ri@setkey\ri@key{\ri@qid#1}\def\ri@lab{#1}%
  \setbox\ri@labbox=\vtop{\hsize=\dimexpr\ri@L-1.6mm\relax\bfseries\raggedright\parskip=0pt
      \strut\ri@lab\ifcsname ri@mk@\ri@key\endcsname\ifx\ri@lab\@empty\else\\\fi
      [\csname ri@mk@\ri@key\endcsname]\fi\strut\par}%
  \xdef\ri@curH{\number\dimexpr\ht\ri@labbox+\dp\ri@labbox\relax}%
  \nointerlineskip\zsavepos{ri-s\the\ri@n}\nointerlineskip\nobreak
  \vbox to 0pt{\hbox{\hskip-\dimexpr\ri@L+\ri@pad-1.2mm\relax\box\ri@labbox}\vss}\nobreak\prevdepth=\dp\strutbox
  \vskip-\parskip\ignorespaces}
% ---- qparts / \pt: question mode records labels (for the marks); solution mode makes rows
\NewCommandCopy\ri@qparts@begin\qparts \let\ri@qparts@end\endqparts
\renewenvironment{qparts}[1][2.6em]{\advance\ri@depth\@ne
  \ifri@sol\else\ri@qparts@begin[#1]\fi}{\ifri@sol\par\else\ri@qparts@end\fi}
\def\ri@qpartsname{qparts}
\renewcommand{\pt}[1]{\ifnum\ri@depth<2 \gdef\ri@outer{#1}\gdef\ri@cur{#1}\else\expandafter\gdef\expandafter\ri@cur\expandafter{\ri@outer#1}\fi
  \ifri@sol\expandafter\ri@solpt\else\expandafter\ri@qpt\fi{#1}}
\def\ri@qpt#1{\expandafter\gdef\csname ri@hp@\ri@qid\endcsname{}\item[#1]}
\def\ri@solpt#1{\@ifnextchar\begin\ri@ptpeek{\ri@row{\ri@cur}}}
\def\ri@ptpeek\begin#1{\def\ri@tmp{#1}\ifx\ri@tmp\ri@qpartsname
   \expandafter\@firstoftwo\else\expandafter\@secondoftwo\fi
  {\begin{qparts}}{\ri@row{\ri@cur}\begin{#1}}}
% marks: record [n] against the current question/part
\let\ri@mk\mk
\renewcommand{\mk}[1]{\ri@setkey\ri@k{\ri@qid\ri@cur}\ri@setkey\ri@mv{#1}\expandafter\xdef\csname ri@mk@\ri@k\endcsname{\ri@mv}\ri@mk{#1}}
% \slabel (Ch10 Test answers): letter labels set the outer part, roman ones are appended;
% consecutive \slabel calls merge into one row label
\def\ri@romans{(i)(ii)(iii)(iv)(v)(vi)(vii)}
\def\ri@firstpart(#1)#2\@nil{(#1)}
\def\slabel#1{\ifri@sol\expandafter\ri@solslabel\else\expandafter\ri@qslabel\fi{#1}}
\def\ri@qslabel#1{\par\smallskip\noindent\textbf{#1}\par\nobreak}
\def\ri@solslabel#1{\@expandtwoargs\in@{#1}{\ri@romans}%
  \ifin@\expandafter\gdef\expandafter\ri@cur\expandafter{\ri@outer#1}\else\gdef\ri@cur{#1}\xdef\ri@outer{\ri@firstpart#1\@nil}\fi
  \@ifnextchar\slabel{}{\ri@row{\ri@cur}}}
% Ch11 spart{label}: a row; soln: the table
\def\spart#1{\gdef\ri@cur{#1}\ri@row{#1}}\def\endspart{\par}
\def\soln{\@ifnextchar[{\ri@solnopt}{\risol}}\def\ri@solnopt[#1]{\risol}\def\endsoln{\endrisol}
% ---- working that started on the label line is left-aligned in the working column
\AtBeginDocument{\@ifundefined{topalign}{}{%
  \renewenvironment{topalign}{\par\noindent\hfill$\displaystyle\begin{aligned}[t]}{\end{aligned}$\hfill\null\par\smallskip}%
  \renewcommand{\topdisplay}[1]{\par\noindent\hfill$\displaystyle #1$\hfill\null\par\smallskip}}}
% ---- Ch11 side-by-side (text | sketch): the sketch shrinks to fit its half of the working column
\AtBeginDocument{\@ifundefined{sbsplain}{}{%
  \renewenvironment{sbsplain}[1]{%
    \sbs@start\begin{minipage}[t]{#1\linewidth}\vspace{0pt}\sbs@shift\raggedright%
    \def\sbsnext{\end{minipage}\hfill\begin{minipage}[t]{\dimexpr0.97\linewidth-#1\linewidth\relax}\vspace{0pt}\sbs@shift\centering
      \begin{adjustbox}{max width=\linewidth}}%
  }{\end{adjustbox}\end{minipage}\par}}}
% ---- figures in the working column shrink to fit it (only if wider), colours kept
\newcommand{\rifitfig}[1]{\expandafter\let\csname ri@orig@\string#1\endcsname#1%
  \def#1{\@ifnextchar\bgroup{\ri@fitarg#1}{\ri@fitnoarg#1}}}
\def\ri@fitarg#1#2{\ifri@sol\adjustbox{max width=\linewidth}{\csname ri@orig@\string#1\endcsname{#2}}\else\csname ri@orig@\string#1\endcsname{#2}\fi}
\def\ri@fitnoarg#1{\ifri@sol\adjustbox{max width=\linewidth}{\csname ri@orig@\string#1\endcsname}\else\csname ri@orig@\string#1\endcsname\fi}
\@ifundefined{solfig}{}{\renewcommand{\solfig}[2][]{\par\begin{center}\adjustbox{max width=\linewidth}{\begin{tikzpicture}\begin{axis}[qfig, width=6.6cm, height=5cm, #1]#2\end{axis}\end{tikzpicture}}\end{center}}}
% ---- question + solution; each question starts on a new page (as RI)
\def\ri@qid{}
\let\ri@putQ\putQ
\renewcommand{\putQ}[1]{\gdef\ri@qid{#1}\global\ri@depth=0 \gdef\ri@cur{}\ri@putQ{#1}}
% \risolQ{n}{body macro name}: question n + its solution body in the table;
% \risolQraw: body already contains its own soln environment (Ch11 Test)
\newcommand{\ri@newq}{\ifri@fresh\global\ri@freshfalse\else\clearpage\fi}
% keep the whole table together if it fits on a page; otherwise make sure at least its
% header and first row fit before starting it
\newcommand{\ri@keep}[1]{\par
  \ifcsname ri@total@#1\endcsname
    % hysteresis: a table kept whole in the previous run stays whole up to textheight-4.5mm,
    % so a total near the limit cannot flip between runs
    \ifdim\csname ri@total@#1\endcsname<\dimexpr\textheight-\ifcsname ri@kept@#1\endcsname4.5\else8\fi mm\relax
      \protected@write\@auxout{}{\string\ri@setkept{#1}}%
      \needspace{\dimexpr\csname ri@total@#1\endcsname+4mm\relax}%
    \else\ifcsname ri@first@#1\endcsname
      \ri@rowh{\csname ri@first@#1\endcsname}%
      \ifdim\ri@h<\dimexpr\textheight-20mm\relax\needspace{\dimexpr\ri@h+18mm\relax}\fi\fi\fi\fi}
\newcommand{\risolQ}[2]{\ri@newq\putQ{#1}\vspace{-3ex}\ri@keep{#1}\begin{risol}\csname #2@#1\endcsname\par\end{risol}}
\newcommand{\risolQraw}[2]{\ri@newq\putQ{#1}\vspace{-3ex}\ri@keep{#1}\csname #2@#1\endcsname\par}
\renewcommand{\putQS}[1]{\risolQ{#1}{diS}}
% end marker only if it fits on the last page (no page holding just the marker)
\renewcommand{\diend}[1][END OF SOLUTIONS]{\par\ifdim\dimexpr\pagegoal-\pagetotal\relax<4\baselineskip\else\vspace{1.5ex}\begin{center}\textbf{#1}\end{center}\fi}
\makeatother
